\everymath{\displaystyle}

%%%%%%%%%%%%%%%%%%%%%%%%%%%%%%%%%%%%%%%%
%           main body
%%%%%%%%%%%%%%%%%%%%%%%%%%%%%%%%%%%%%%%%

%-----------------------------------
%	TITLE PAGE
%-----------------------------------

\title[代数封闭,\\ 分裂域及正规扩张]{\texorpdfstring{\LARGE \bfseries 第11部分\quad 代数封闭, 分裂域及正规扩张}{第11部分 代数封闭, 分裂域及正规扩张}}
\author[抽象代数 2]{\Large \bfseries 抽象代数 2}
\date{}

%------------------------------------------------

\begin{document}

%------------------------------------------------

\begin{frame}

\titlepage % Print the title page as the first slide

\end{frame}

%------------------------------------------------

\section[分裂域]{多项式集的分裂域}

%------------------------------------------------

\begin{frame}{多项式集与分裂域}

$F$ 为一个域, $F[x]$ 为多项式环, 令 $S \subset F[x]$,
$S$ 中元素为正次数的多项式, $A$ 为 $S$ 中多项式的零点的集合.

\pause

\begin{Def}[分裂域]
一个代数扩张 $E/F$ 称为 $F$ 上的多项式集 $S$ 的一个分裂域,
若 $A \subset E$, 而且 $E = F(A)$.
\end{Def}

\end{frame}

%------------------------------------------------

\begin{frame}{分裂域的存在性}

\begin{itemize}
\item \circled{1} $S = \{ f(x) \}$:\\
  由不可约多项式的根引理, 至少有 $K/F$, 使得 $K$ 中包含 $f(x)$ 的一个零点
  $\alpha_1$, 那么 $f(x) = (x - \alpha_1) f_1(x)$.
  对 $f_1(x)$ 也可进行同样讨论,
  故 $\exists$ 域扩张 $L/F$, $[L:F] < \infty$, 且 $L$ 中包含 $f(x)$ 的全部零点
  $A$, 令 $E = F(A)$, 即有 $F \subset F(A) \subset L$, $E$ 为一个分裂域;
\item \circled{2} $S = \{ f_1(x), f_2(x), \dots, f_n(x) \}$:\\
  令 $f(x) = \prod_{i=1}^{n} f_i(x)$ 即归结到 \circled{1};
\item \circled{3} $S = F[x]$.
\end{itemize}

\end{frame}

%------------------------------------------------

\section[闭包]{代数闭域与代数闭包}

%------------------------------------------------

\begin{frame}{代数闭域与代数闭包}

\begin{Def}[代数闭域]
一个域 $E$ 为一个代数闭域, 若 $E$ 上任一不可约多项式为一次多项式.
\end{Def}

\pause

若 $E/F$ 为域扩张, 称 $E$ 为 $F$ 的一个代数闭包, 若:
\begin{itemize}
\item[(i)] $E/F$ 为代数扩张;
\item[(ii)] $E$ 为代数闭域.
\end{itemize}
以 $\bar{F}$ 记 $F$ 的代数闭包.

\pause

注: $E$ 为代数闭域 $\iff$ 任意代数扩张 $L/E$, 有 $L = E$.

\end{frame}

%------------------------------------------------

\begin{frame}{例: 代数闭域}

\begin{eg}
$F = \mathbb{C}$ 为一个代数闭域,
$\bar{F} = \bar{\mathbb{Q}}_c = \{ z \in \mathbb{C} \mid z \text{ 为代数数} \}$
为代数闭域.
令 $F = \mathbb{Q}$, 有 $F$ 的代数闭包 $\bar{F} = \bar{\mathbb{Q}}_c$.
\end{eg}

\end{frame}

%------------------------------------------------

\begin{frame}{代数闭包的存在性}

\begin{prop}[代数闭包的存在性]
任一个域 $F$ 的代数闭包 $\bar{F}$ 存在.
\end{prop}

\startproof[bg=yellow, fg=red](证明)
令 $T = \{ E \mid E/F \text{ 为代数扩张} \}$.
$T$ 非空, 因为至少有 $F \in T$. 在包含关系下, $T$ 为非空的偏序集.
对任一全序子集 $E_1 \subset E_2 \subset \dots \subset E_n \subset \dots$,
令 $E = \bigcup_{i=1}^{\infty} E_i$, 则 $E/F$ 为代数扩张, 即 $E \in T$,
$\{ E_i \}$ 在 $T$ 中有上界.
由 Zorn 引理, $T$ 中有极大元素, 令其为 $E$. $E$ 即为 $F$ 的一个代数闭包.
\qed

\end{frame}

%------------------------------------------------

\begin{frame}{分裂域的存在唯一性}

\begin{thm}[分裂域的存在唯一性]
$F$ 为一个域, $S \subset F[x]$, $S$ 的分裂域存在,
而且 $S$ 的任意两个分裂域都是 $F$-同构的.
\end{thm}

\startproof[bg=yellow, fg=red](证明)
令 $A$ 为 $S \subset F[x]$ 的零点集, $\bar{F}$ 存在, 而 $A \subset \bar{F}$,
那么 $F \subset F(A) \subset \bar{F}$, 即 $F(A)$ 应是 $S$ 在 $F$ 上的分裂域.
\qed

\end{frame}

%------------------------------------------------

\section[开拓]{分裂域的开拓}

%------------------------------------------------

\begin{frame}{分裂域的开拓}

\begin{lem}[分裂域的开拓]
$F \xrightarrow{\ \sigma\ } F_1$ 为域同构, $S \subset F[x]$,
令 $S' = \{ \sigma(f(x)) \mid f(x) \in S \}$ 为 $F_1$ 上 $S$ 所对应的多项式集.
令 $E$ 为 $S$ 的一个分裂域, $E'$ 为 $S'$ 在 $F_1$ 上的一个分裂域,
那么有域同构 $E \xrightarrow{\ \bar{\tau}\ } E'$, 满足 $\bar{\tau}|_F = \sigma$.
\end{lem}

\end{frame}

%------------------------------------------------

\begin{frame}{分裂域的开拓的证明 (一)}

\startproof[bg=yellow, fg=red](证明)
令 $A$ 为 $S$ 的零点集, $A'$ 为 $S'$ 的零点集, $E = F(A)$, $E' = F_1(A')$.
考虑集合
\[
T = \{ (L, \varphi) \mid F \subset L \subset E,\
L \xrightarrow{\ \varphi\ } E' \text{ 为同态},\ \varphi|_F = \sigma \}.
\]
任一 $\varphi$ 为单同态, $L \cong \varphi(L) \subset E'$.
$T$ 非空, 因为 $(F, \sigma) \in T$, 定义一个偏序:
$(L_1, \varphi_1) \preccurlyeq (L_2, \varphi_2)$,
若 $L_1 \subset L_2$, 且 $\varphi_2|_{F_1} = \varphi_1$.
那么 $T$ 为一个非空偏序集合.
\[
\text{若 } (L_1, \varphi_1) \preccurlyeq (L_2, \varphi_2) \preccurlyeq \dots
\preccurlyeq (L_n, \varphi_n) \preccurlyeq \dots
\]
令 $L = \bigcup_{i=1}^{\infty} L_i$, 有 $F \subset L \subset E$;
$\forall\ \alpha \in L$, 存在 $L_i$, 使得 $\alpha \in L_i$,
令 $\varphi(\alpha) = \varphi_i(\alpha)$.
易知 $\varphi$ 为 $L \to E'$ 的一个同态, $\varphi|_F = \sigma$, 即 $(L, \varphi) \in T$,
即 $\{ (L_i, \varphi_i) \}$ 在 $T$ 中有上界, 由 Zorn 引理,
$T$ 中有极大元素, 令其为 $(T, \bar{\tau})$, 可以证明, $L = E$, $\bar{\tau}(E) = E'$.

\end{frame}

%------------------------------------------------

\begin{frame}{分裂域的开拓的证明 (二)}

\startproof[bg=yellow, fg=red](证明, 续)
(i) 若 $L \neq E$, 由 $E = F(A)$ 知 $\exists\ \alpha \in A$, 使得 $\alpha \notin L$,
显见, $F \subset L$, $A \subset L$, 那么 $F(A) \subset L$.
由同构的开拓引理, $L \cong E'$, $\alpha \notin L$,
$\alpha$ 对应的 $S'$ 中的零点 $\beta$, $\bar{\tau}$ 开拓为
$L(\alpha) \xrightarrow{\ \tau'\ } E'$, 满足 $\tau'|_L = \bar{\tau}$,
则 $(L, \bar{\tau}) \preccurlyeq (L(\alpha), \tau')$,
但 $(L, \bar{\tau}) \neq (L(\alpha), \tau')$,
这与 $(L, \bar{\tau})$ 的极大性矛盾. 故必有 $L = E$.
\[
\text{(ii) } \bar{\tau}(E) = \bar{\tau}(F(A)) = \sigma(F)(\bar{\tau}(A)) = F_1(A') = E'.
\]
综上 有 $E \xrightarrow{\ \bar{\tau}\ } E'$ 为域同构, 且 $\bar{\tau}|_F = \sigma$.
即: 域同构 $F \xrightarrow{\ \sigma\ } F_1$ 可以开拓到对应的分裂域上.
\qed

\end{frame}

%------------------------------------------------

\begin{frame}{分裂域的推论}

\begin{cor}[分裂域的推论]
\begin{itemize}
\item[(i)] $F$ 上任一多项式集 $S$ 的任意两个分裂域都是 $F$-同构的.
  特别地, $F$ 上任意两个代数闭包都是 $F$-同构的;
\item[(ii)] 若 $F \xrightarrow{\ \sigma\ } F'$ 为域同构, 则 $\sigma$ 可以开拓到
  $\bar{F} \xrightarrow{\ \bar{\tau}\ } \bar{F}'$, $\bar{\tau}|_F = \sigma$;
\item[(iii)] $E/F$ 为代数扩张, $\bar{F}$ 为 $F$ 的代数闭包, 存在 $F$-同态
  $\sigma: E \xrightarrow{\ \tau\ } \bar{F}$, 即 $E$ 可以嵌入到 $\bar{F}$.
\end{itemize}
\end{cor}

\startproof[bg=yellow, fg=red](证明)
(i) 取分裂域的开拓引理中的 $F_1 = F$, $\sigma = 1$, 有 $S' = S$,
由分裂域的开拓引理, 即有命题 (i) 成立
(注: 在同构意义下, $\bar{F}$ 有意义);
(ii) 在分裂域的开拓引理中, 取 $S = F[x]$.
\qed

\end{frame}

%------------------------------------------------

\begin{frame}{分裂域的推论的证明 (iii)}

\startproof[bg=yellow, fg=red]((iii) 的证明)
若 $E/F$ 为代数扩张, 那么 $\bar{E}$ 也是 $F$ 的代数闭包,
因为 $\bar{E}/E$, $E/F$ 是代数的, 则 $\bar{E}/F$ 也是代数的, 故 $\bar{E} = \bar{F}$.
由 (iii) 知 $F$ 上两个代数闭包 $\bar{E}$ 及 $\bar{F}$ 为 $F$-同构的.
令 $\bar{E} \xrightarrow{\ \bar{\tau}\ } \bar{F}$, $\bar{\tau}|_F = 1$,
那么 $\bar{\tau}|_E: E \to \bar{F}$ 为一个 $F$-同态, 或一个 $F$-嵌入.
\qed

\end{frame}

%------------------------------------------------

\section[正规]{正规扩张}

%------------------------------------------------

\begin{frame}{正规扩张}

\begin{Def}[正规扩张]
一个代数扩张 $E/F$ 称为 $F$ 上的一个正规扩张
(normal extension), 若 $F$ 上任一不可约多项式 $p(x)$ 在 $E$ 中有一个零点,
则 $E$ 中包含 $p(x)$ 的全部零点.
\end{Def}

\end{frame}

%------------------------------------------------

\begin{frame}{正规扩张的刻画}

\begin{lem}[正规扩张的刻画]
$E/F$ 为代数扩张, 则下述命题等价:
\begin{itemize}
\item[(i)] $E$ 为 $F$ 上某一个多项式集 $S$ 的分裂域;
\item[(ii)] $E \xrightarrow{\ \bar{\tau}\ } \bar{F}$ 为 $F$-同态,
  则 $\bar{\tau}(E) = E$, 即 $\bar{\tau} \in \mathrm{Gal}(E/F)$;
\item[(iii)] $F \subset L \subset E$, $L$ 为任一中间域,
  $L \xrightarrow{\ \sigma\ } \bar{F}$ 为 $F$-同态, 则存在
  $\bar{\tau} \in \mathrm{Gal}(E/F)$, 且 $\bar{\tau}|_L = \sigma$;
\item[(iv)] $E/F$ 为正规扩张.
\end{itemize}
\end{lem}

\end{frame}

%------------------------------------------------

\begin{frame}{正规扩张的刻画的证明 (一)}

\startproof[bg=yellow, fg=red]((i) $\Longrightarrow$ (ii) 的证明)
令 $E$ 为 $F$ 上多项式集 $S$ 的分裂域, $A$ 为 $S$ 的零点集, 即 $E = F(A)$.
任取一个 $F$-同态 $E \xrightarrow{\ \bar{\tau}\ } \bar{F}$,
\[
\bar{\tau}(E) = \bar{\tau}(F(A)) = F(\bar{\tau}(A)) = F(A) = E.
\]

\end{frame}

%------------------------------------------------

\begin{frame}{正规扩张的刻画的证明 (二)}

\startproof[bg=yellow, fg=red]((ii) $\Longrightarrow$ (iii) 的证明)
若 $F \subset L \subset E$ 为中间域, $L \xrightarrow{\ \sigma\ } \bar{F}$
为 $F$-同态, 令 $\sigma(L) = L'$,
即 $F \subset L \subset E \subset \bar{F}$.
$L \xrightarrow{\ \cong\ } L'$ 开拓到代数闭包上,
并注意到 $\bar{F}$ 为 $L$ 及 $L'$ 的代数闭包,
即 $\exists\ \bar{\tau}: \bar{F} \to \bar{F}$ 为 $F$-同构, 且 $\bar{\tau}|_L = \sigma$.
由 (ii) $\bar{\tau}|_E = \bar{\tau}$, $\bar{\tau}(E) = E$,
$\sigma(L) = \bar{\tau}(L) \subset E$ 且 $\bar{\tau}|_L = \sigma$.

\end{frame}

%------------------------------------------------

\begin{frame}{正规扩张的刻画的证明 (三)}

\startproof[bg=yellow, fg=red]((iii) $\Longrightarrow$ (iv) 的证明)
$p(x)$ 为 $F$ 上不可约多项式, 在 $E$ 中有一个 0 点 $\alpha \in E$,
若 $\beta$ 为 $p(x)$ 的另一个零点, 由同构的开拓引理,
$\sigma = 1_F$ 开拓到 $F(\alpha) \cong F(\beta)$,
满足 $\bar{\tau}|_F = 1_F$, 注意到 $F(\alpha) \subset E$,
由 (iii) 知 $\bar{\tau}(F(\alpha)) \subset E$,
而 $\bar{\tau}(F(\alpha)) = F(\beta)$, 那么 $\beta \in E$,
从而 $E/F$ 为正规扩张.

\end{frame}

%------------------------------------------------

\begin{frame}{正规扩张的刻画的证明 (四)}

\startproof[bg=yellow, fg=red]((iv) $\Longrightarrow$ (i) 的证明)
若 $E/F$ 为正规扩张, 令 $S = \{ \text{在 } E \text{ 中有一个零点的 }
F \text{ 上的多项式 } p(x) \}$,
令 $S$ 的零点集为 $A$.
因为 $E/F$ 为正规扩张, 故 $F(A) \subset E$.
反过来, $\forall\ \alpha \in E$, $\alpha$ 的极小多项式 $p(x) \in S$,
那么 $\alpha \in A$, 故 $E \subset F(A)$.
综上知 $E = F(A)$.
\qed

\end{frame}

%------------------------------------------------

\begin{frame}{有限正规扩张与分裂域}

\begin{cor}[有限正规扩张与分裂域]
任一有限的正规扩张 $E/F$ $\iff$ $E$ 为 $F$ 上的一个多项式 $p(x)$ 的分裂域.
\end{cor}

\startproof[bg=yellow, fg=red](证明)
$E$ 为有限正规扩张:
若 $E$ 为 $p(x)$ 的分裂域, 令 $E = F(\alpha_1, \dots, \alpha_n)$,
$\alpha_1, \dots, \alpha_n$ 为 $p(x)$ 全部零点,
由定义即知 $E/F$ 为正规扩张, 且 $[E:F] \leq (\deg p(x))!$.
反过来, 若 $[E:F] \leq n$, 且为正规扩张,
令 $E = F(\alpha_1, \dots, \alpha_m)$, 其中 $\alpha_1, \dots, \alpha_m$
为 $F$ 上的代数元素,
令 $g_i(x)$ 为 $\alpha_i$ 在 $F$ 上的极小多项式,
令 $f(x) = g_1(x) \dotsm g_m(x)$, 那么 $E$ 为 $f(x)$ 的分裂域.
\qed

\end{frame}

%------------------------------------------------

\section[闭包2]{正规闭包}

%------------------------------------------------

\begin{frame}{正规闭包}

正规闭包: 若 $E/F$ 为代数扩张, $F \subset E \subset N$, $N/F$ 为正规扩张,
且 $\forall\ L$, $F \subset E \subset L \subset N$,
满足 $L/F$ 为正规扩张 $\Longrightarrow L = N$,
此时 $N$ 称为 $E/F$ 的正规闭包.

\end{frame}

%------------------------------------------------

\begin{frame}{正规闭包引理}

\begin{lem}[正规闭包]
\begin{itemize}
\item[(i)] 任意代数扩张 $E/F$ 的正规闭包 $N/F$ 存在,
  且在 $F$-同构下是唯一存在的;
\item[(ii)] 任意有限扩张 $E/F$ 的正规闭包 $N/F$ 为有限扩张.
\end{itemize}
\end{lem}

\startproof[bg=yellow, fg=red]((i) 的证明)
若 $E/F$ 为代数扩张, 令
$S = \{ g_{\alpha}(x) \mid \alpha \in E,\ g_{\alpha}(x) \text{ 为 } \alpha \text{ 的极小多项式} \}$,
令 $N$ 为 $S$ 的分裂域, 即有 $F \subset E \subset N$, 且 $N/F$ 为正规扩张.
现任取 $L$, $F \subset E \subset L \subset N$.
若 $L$ 为正规扩张, 由于 $E \subset L$, 故 $S$ 的零点集 $A \subset L$,
那么 $N = F(A) \subset L$, 即 $L = N$.
$\Longrightarrow N/F$ 为 $E/F$ 的正规闭包.
\qed

\end{frame}

%------------------------------------------------

\begin{frame}{正规闭包引理的证明 (ii)}

\startproof[bg=yellow, fg=red]((ii) 的证明)
$[E:F] = n$, $E = F(\alpha_1, \dots, \alpha_m)$,
令 $g_i(x)$ 为 $\alpha_i$ 的极小多项式,
令 $g(x) = \prod_i g_i(x)$, 令 $N$ 为 $g(x)$ 的分裂域,
那么 $N/F$ 为 $E/F$ 的正规闭包, 且 $[N:F] \leq (\deg g(x))!$,
所以 $N/F$ 也是有限扩张.
\qed

\end{frame}

%------------------------------------------------

\section[有限域]{有限域的构造}

%------------------------------------------------

\begin{frame}{有限域的构造}

\begin{eg}[有限域的构造]
$F$ 为有限域, 若 $|F| < \infty$, $F$ 又称为 Galois 域.
若 $F$ 为有限域, 则 $\chi(F) = p$ 为一个素数, 那么 $F_p \subset F$,
令 $[F:F_p] = n$, 那么 $|F| = p^n$.
\end{eg}

\begin{itemize}
\item 结论: 任一有限域 $F \Longrightarrow |F|$ 为某个素数的方幂;
\item 反过来, 任意给定一个素数 $p$, 及 $n \in \mathbb{N}$,
  唯一存在一个有限域 $F$, 满足 $|F| = p^n$.
\end{itemize}

\end{frame}

%------------------------------------------------

\begin{frame}{结论的证明}

\startproof[bg=yellow, fg=red](证明)
$F_p$ 为 $p$ 元有限域, $\bar{F}_p$ 为 $F_p$ 的代数闭包.
$g(x) = x^{p^n} - x \in F_p[x]$, 令 $F$ 为 $g(x)$ 的分裂域,
$F_p \subset F \subset \bar{F}_p$.
由于 $g(x)$ 无重根, 令 $A = \{ \alpha_1, \dots, \alpha_{p^n} \} \subset F$,
为 $g(x)$ 的根的集合. 易证明 $A$ 为一个域.
又因为 $\forall\ \alpha \in F_p$, 有 $\alpha^p - \alpha = 0$,
故 $F_p \subset A \subset \bar{F}_p$,
那么 $A = F_p(A) = F$, $|F| = |A| = p^n$.
又有对任意 $|F| = p^n$, $F$ 为 $x^{p^n} - x$ 的分裂域,
即两个有 $p^n$ 元的有限域同构.
\qed

\end{frame}

%------------------------------------------------

\begin{frame}{有限域的 Galois 扩张}

\begin{prop}[有限域的 Galois 扩张]
任一 $p^n$ 元的有限域 $F$ 为 $F_p$ 上的 $n$ 次 Galois 扩张.
\end{prop}

\startproof[bg=yellow, fg=red](证明)
$[F:F_p] = n$, 令 $\sigma \in \mathrm{Gal}(F/F_p)$,
定义为, $\forall\ \alpha \in F$, $\sigma(\alpha) = \alpha^p$,
易验证知 $\sigma \in \mathrm{Gal}(F/F_p)$, $\sigma|_{F_p} = 1_{F_p}$.
又因为 $\forall\ \alpha \in F$, $\alpha^{p^n} = \alpha$, 故 $\sigma^n = 1$
\[
\Longrightarrow \mathrm{Gal}(F/F_p) = \langle \sigma \rangle
= \{ \sigma, \sigma^2, \dots, \sigma^n \}.
\]
$\sigma$ 称为 $F$ 上的 Frobenius 自同构.
\qed

\end{frame}

%------------------------------------------------

\begin{frame}{有限域的正规性与 Galois 性}

\begin{itemize}
\item 结论: 任一有限域 $F$, $|F| = p^n \Longrightarrow F/F_p$ 为正规扩张,
  也是 Galois 扩张.
\end{itemize}

\pause

\begin{eg}
令 $F$ 为有限域, $|F| = p^n$, $F(x)$ 为单变量的有理函数域,
$F \subset F(x^p) \subset F(x)$, 那么 $F(x)/F(x^p)$ 为 $p$ 次正规扩张,
但非 Galois 扩张.
\end{eg}

\startproof[bg=yellow, fg=red](证明)
$F(x) = F(x^p)(x)$, $x$ 为 $F(x^p)$ 上的代数元素,
$g(t) = t^p - x^p \in F(x^p)[t]$ 为 $x$ 的极小多项式.
由于 $x$ 为 $g(t)$ 的全部零点, 故 $F(x)$ 为 $g(t)$ 在 $F(x^p)$ 上的分裂域,
即知 $F(x)/F(x^p)$ 为 $p$ 次正规扩张.
又已证明 $\mathrm{Gal}(F(x)/F(x^p)) = 1$,
故 $F(x)/F(x^p)$ 不是 Galois 扩张.
推广到 $F(x_1, \dots, x_n)/F(x_1^p, \dots, x_n^p)$ 为 $p^n$ 次正规扩张,
但其 Galois 群为平凡的, 因此非 Galois 扩张.
\qed

\end{frame}

%------------------------------------------------

\end{document}
